% ============================================================
%  Raport: Bamboo Garden Trimming – Gry Kombinatoryczne
%  Politechnika Warszawska, 2024/2025
% ============================================================
\documentclass[12pt,a4paper]{article}

% --- Kodowanie, język i fonty -----------------------------------------
\usepackage{iftex}
\ifPDFTeX
  \usepackage[T1]{fontenc}
  \usepackage[utf8]{inputenc}
  \usepackage{lmodern}
\else
  \usepackage{fontspec}
  \setmainfont{lmroman12-regular.otf}[
    BoldFont = lmroman12-bold.otf,
    ItalicFont = lmroman12-italic.otf,
    BoldItalicFont = lmroman10-bolditalic.otf
  ]
\fi
\usepackage[polish]{babel}

% --- Strona i typografia ----------------------------------------------
\usepackage[a4paper, top=2.5cm, bottom=2.5cm, left=3cm, right=2.5cm]{geometry}
\usepackage{microtype}
\usepackage{setspace}
\onehalfspacing
\usepackage{parskip}
\setlength{\parindent}{0pt}
\emergencystretch=2em

% --- Matematyka --------------------------------------------------------
\usepackage{amsmath, amssymb, amsthm}
\usepackage{mathtools}

% --- Twierdzenia -------------------------------------------------------
\theoremstyle{plain}
\newtheorem{twierdzenie}{Twierdzenie}[section]
\newtheorem{lemat}[twierdzenie]{Lemat}
\newtheorem{wniosek}[twierdzenie]{Wniosek}

\theoremstyle{definition}
\newtheorem{definicja}[twierdzenie]{Definicja}
\newtheorem{przyklad}[twierdzenie]{Przykład}

\theoremstyle{remark}
\newtheorem{uwaga}[twierdzenie]{Uwaga}

% --- Algorytmy ---------------------------------------------------------
\usepackage{algorithm}
\usepackage{algpseudocode}
\floatname{algorithm}{Algorytm}
\renewcommand{\algorithmicrequire}{\textbf{Wejście:}}
\renewcommand{\algorithmicensure}{\textbf{Wyjście:}}

% --- Tabele i grafika --------------------------------------------------
\usepackage{booktabs}
\usepackage{tabularx}
\usepackage{array}
\newcolumntype{Y}{>{\raggedright\arraybackslash}X}
\usepackage{graphicx}
\usepackage{tikz}
\usetikzlibrary{positioning, arrows.meta, decorations.pathreplacing, calc}
\usepackage{pgfplots}
\pgfplotsset{compat=1.18}
\usepackage{xcolor}
\definecolor{bgreen}{HTML}{2d8a3e}
\definecolor{bamber}{HTML}{b45309}
\definecolor{bblue}{HTML}{1d4ed8}
\definecolor{bred}{HTML}{dc2626}
\definecolor{bgray}{HTML}{f4f0e8}

% --- Inne pakiety ------------------------------------------------------
\usepackage{enumitem}
\usepackage{titlesec}
\usepackage{fancyhdr}
\usepackage{hyperref}
\usepackage{cleveref}
\hypersetup{
  colorlinks=true,
  linkcolor=bblue,
  citecolor=bamber,
  urlcolor=bblue
}
\crefname{section}{sekcji}{sekcji}
\Crefname{section}{Sekcja}{Sekcje}
\crefname{twierdzenie}{twierdzenia}{twierdzeń}
\Crefname{twierdzenie}{Twierdzenie}{Twierdzenia}
\crefname{definicja}{definicji}{definicji}
\Crefname{definicja}{Definicja}{Definicje}
\crefname{wniosek}{wniosku}{wniosków}
\Crefname{wniosek}{Wniosek}{Wnioski}

% --- Styl nagłówków ---------------------------------------------------
\titleformat{\section}{\large\bfseries}{\thesection.}{0.6em}{}[\titlerule]
\titleformat{\subsection}{\normalsize\bfseries}{\thesubsection.}{0.5em}{}

% --- Nagłówek i stopka ------------------------------------------------
\pagestyle{fancy}
\fancyhf{}
\fancyhead[L]{\small\textit{Bamboo Garden Trimming}}
\fancyhead[R]{\small\textit{Gry Kombinatoryczne – PW}}
\fancyfoot[C]{\small\thepage}
\renewcommand{\headrulewidth}{0.4pt}

% --- Makra matematyczne -----------------------------------------------
\newcommand{\hmax}{h_{\max}}
\newcommand{\makespan}{\mathcal{M}}
\newcommand{\RF}{\textsc{Reduce-Fastest}}
\newcommand{\RM}{\textsc{Reduce-Max}}
\newcommand{\OPT}{\textsc{Opt}}

% ======================================================================
\begin{document}

% ======= STRONA TYTUŁOWA ==============================================
\begin{titlepage}
  \centering
  \vspace*{1.5cm}

  {\Large\bfseries Politechnika Warszawska\\[0.3em]
  Wydział Matematyki i Nauk Informacyjnych}

  \vspace{1.2cm}
  \hrule height 1.5pt
  \vspace{0.6cm}

  {\LARGE\bfseries Bamboo Garden Trimming}\\[0.5em]
  {\large Przygotowywanie ogrodu bambusowego –\\
  algorytmiczne aspekty problemu harmonogramowania wieczystego}

  \vspace{0.6cm}
  \hrule height 1.5pt
  \vspace{1.5cm}

  {\large Projekt zaliczeniowy z przedmiotu}\\[0.3em]
  {\large\bfseries Gry Kombinatoryczne}

  \vspace{2cm}

  \begin{tabular}{ll}
    \textbf{Autorzy:}   & Michał Bober \\
                        & Piotr Komisarczyk \\
                        & Igor Staręga \\[0.8em]
    \textbf{Data:}      & \today
  \end{tabular}

  \vfill

  \begin{center}
    \small Praca oparta na artykule:\\
    D.~Bilò, L.~Gualà, S.~Leucci, G.~Proietti, G.~Scornavacca,\\
    \textit{Cutting bamboo down to size},\\
    Theoretical Computer Science 909 (2022), ss.~54–67.
  \end{center}
\end{titlepage}

% ======= STRESZCZENIE =================================================
\begin{abstract}
\noindent
Niniejszy raport omawia problem \emph{Bamboo Garden Trimming} (BGT) –
harmonogramowania wieczystego sekwencji cięć robota-ogrodnika dbającego
o ogród bambusowy. Zadanie polega na wyznaczeniu nieskończonego harmonogramu
cięć minimalizującego \emph{makespan}, czyli największą wysokość kiedykolwiek
osiągniętą przez dowolny bambus. Przedstawiamy formalny model BGT,
intuicję stojącą za redukcją do problemu Pinwheel Scheduling oraz dwie
naturalne strategie: \RM{} (przycinaj aktualnie najwyższy bambus) i
\RF$(x)$ (przycinaj najszybciej rosnący bambus spośród tych, które osiągnęły
próg~$x$). Omawiamy najważniejsze wyniki z~pracy Bilò i~wsp.: pierwsze
stałe ograniczenie górne~$9$ dla \RM{} oraz ograniczenie
$\tfrac{3+\sqrt{5}}{2} < 2{,}62$ dla~\RF$(x)$ przy
$x = 1 + \tfrac{1}{\sqrt{5}}$. Ostatnia część raportu opisuje naszą
implementację w~Pythonie i~Streamlit: interaktywną aplikację pozwalającą
symulować, porównywać i~testować omawiane strategie.
\end{abstract}

% ======= SPIS TREŚCI ==================================================
\newpage

\tableofcontents

\newpage

% ======================================================================
\section{Wprowadzenie i motywacja}
% ======================================================================

Problem \emph{Bamboo Garden Trimming} (BGT) jest matematyczną formalizacją
pozornie banalnego zadania utrzymania ogrodu bambusowego. Wyobraźmy sobie
dom nad jeziorem, przed którym rośnie ogród złożony z $n$~bambusów.
Każdy bambus~$b_i$ rośnie o stałą wartość~$h_i > 0$ na dobę.
Robot-ogrodnik może raz dziennie skosić \emph{co najwyżej jeden} bambus,
natychmiast resetując jego wysokość do zera. Pytanie brzmi: jak zaplanować
sekwencję cięć, aby widok na jezioro był jak najmniej zasłonięty?

Problem ten, sformalizowany przez G\k{a}sieńca i współautorów
w~\cite{gasieniec2017}, łączy w sobie zagadnienia z teorii szeregowania zadań,
gier kombinatorycznych oraz struktury danych. Pomimo prostego opisu,
prowadzi do interesujących pytań dotyczących granicy między scenariuszem
statycznym a~adwersarialnym.

\subsection{Znaczenie praktyczne}

Choć historia z~pandą-ogrodnikiem jest celowo lekka, model BGT opisuje
realny typ problemów: ograniczony zasób obsługi musi cyklicznie reagować
na obiekty narastające z~różnymi szybkościami. Problem BGT i~jego
uogólnienia pojawiają się między innymi w~następujących zastosowaniach:
\begin{itemize}[noitemsep]
  \item \textbf{Zarządzanie buforami sieciowymi} – pakiety o różnych
        priorytetach wymagają cyklicznego przetwarzania~\cite{goldwasser2010};
  \item \textbf{Monitorowanie usług w~chmurze} – maszyny wirtualne wymagają
        okresowej inspekcji pod kątem anomalii~\cite{alshamrani2015};
  \item \textbf{Harmonogramowanie czasu rzeczywistego} – zadania z różnymi
        częstotliwościami obsługi muszą być realizowane bez przekroczeń
        terminów~\cite{bender2020};
  \item \textbf{Problem kubka} (\emph{cup game}) – uogólnienie BGT, w którym
        dzienne cięcie może przekroczyć $1$ jednostkę.
\end{itemize}

% ======================================================================
\section{Formalna definicja problemu}
% ======================================================================

\begin{definicja}[Problem BGT]
Ogród zawiera $n$ bambusów $b_1, \ldots, b_n$, gdzie bambus $b_i$ ma
dzienny współczynnik wzrostu $h_i > 0$. Zakładamy, że
\[
  \sum_{i=1}^{n} h_i = 1
\]
(bez straty ogólności – skalowanie nie zmienia struktury optymalnego
harmonogramu). Początkowo wszystkie bambusy mają wysokość $0$. Na końcu
każdego dnia $d \geq 1$ ogrodnik może przyciąć \emph{co najwyżej jeden}
bambus, natychmiast resetując jego wysokość do $0$.

Wysokość bambusu $b_i$ na końcu dnia $d$ (przed decyzją o cięciu) wynosi:
\[
  H_i(d) \;=\; (d - d') \cdot h_i,
\]
gdzie $d' < d$ to ostatni dzień, w którym $b_i$ był cięty
(przyjmujemy $d' = 0$, jeśli $b_i$ nigdy nie był cięty).

Celem jest wyznaczenie \emph{wieczystego harmonogramu} (nieskończonej
sekwencji decyzji o cięciu) minimalizującego \emph{makespan}:
\[
  \makespan \;=\; \max_{i,\,d} H_i(d).
\]
\end{definicja}

\subsection{Normalizacja i założenia}

Warunek $\sum h_i = 1$ jest uzasadniony następująco: jeśli oryginalne
współczynniki sumują się do $S \neq 1$, wystarczy podzielić każde $h_i$
przez $S$. Makespan skaluje się proporcjonalnie, a~struktura optymalnego
harmonogramu pozostaje niezmieniona.

Przyjmujemy standardowe uporządkowanie $h_1 \geq h_2 \geq \cdots \geq h_n > 0$.

\subsection{Ograniczenia na makespan}

Poniższe twierdzenie ustala dolne i górne granice optymalnego makespanu.

\begin{twierdzenie}[Ograniczenia na makespan]\label{thm:bounds}
Dla każdej instancji BGT:
\begin{enumerate}[label=(\roman*), noitemsep]
  \item $\makespan \geq 1$ dla każdego harmonogramu;
  \item istnieje instancja, dla której $\makespan$ może być dowolnie bliski $2$;
  \item makespan $2$ jest osiągalny dla każdej instancji.
\end{enumerate}
\end{twierdzenie}

\begin{proof}
\textbf{(i)} W~każdym dniu suma wysokości wszystkich bambusów rośnie
o~$1$. Jeżeli makespan byłby mniejszy niż $1$, każde pojedyncze cięcie
usuwałoby mniej niż $1$ jednostkę wysokości. W~długim horyzoncie średnia
wysokość usuwana przez ogrodnika byłaby więc mniejsza niż dzienny przyrost,
co prowadziłoby do nieograniczonego wzrostu łącznej wysokości. Sprzeczność.

\textbf{(ii)} Rozważmy dwa bambusy: $h_1 = 1-\varepsilon$ oraz
$h_2 = \varepsilon$. Aby bambus $b_2$ nie rósł bez ograniczeń, musi zostać
od czasu do czasu przycięty. W~dniu, w~którym ogrodnik tnie $b_2$, nie może
jednocześnie przyciąć $b_1$. Nawet jeśli $b_1$ był cięty dzień wcześniej,
to przed kolejną możliwą interwencją urośnie przez dwa dni, osiągając
wysokość $2(1-\varepsilon) = 2-2\varepsilon$. Dla $\varepsilon \to 0$
otrzymujemy rodzinę instancji wymagających makespanu dowolnie bliskiego~$2$.

\textbf{(iii)} Osiągalność makespanu~$2$ wynika z~redukcji do problemu
Pinwheel Scheduling opisanej w~\cref{sec:pinwheel}.
\end{proof}

\begin{figure}[htbp]
\centering
\begin{tikzpicture}[scale=0.9]
  % Oś czasu
  \draw[->] (0,0) -- (8.5,0) node[right] {$d$ (dzień)};
  \draw[->] (0,0) -- (0,4.5) node[above] {$H_i(d)$};

  % Etykiety osi y
  \foreach \y/\label in {1/1, 2/2, 3/3, 4/4}
    \draw (0.05,\y) -- (-0.05,\y) node[left] {\small$\label$};

  % Bambus b1 (h=0.7)
  \draw[bgreen, thick] (0,0) -- (1,0.7) -- (1,0) -- (2,0.7) -- (2,0)
    -- (3,0.7) -- (3,0) -- (4,0.7) -- (4,0) -- (5,0.7) -- (5,0)
    -- (6,0.7) -- (7,1.4) -- (7,0) -- (8,0.7);
  \node[bgreen, right] at (8,0.7) {\small $b_1$};

  % Bambus b2 (h=0.2)
  \draw[bamber, thick] (0,0) -- (5,1.0) -- (5,0) -- (8,0.6);
  \node[bamber, right] at (8,0.6) {\small $b_2$};

  % Bambus b3 (h=0.1)
  \draw[bblue, thick] (0,0) -- (8,0.8);
  \node[bblue, right] at (8,0.8) {\small $b_3$};

  % Linia makespanu
  \draw[bred, dashed, thick] (0,1.4) -- (8,1.4)
    node[right] {\small makespan};

  % Etykiety osi x
  \foreach \x in {1,...,8}
    \draw (\x,0.05) -- (\x,-0.05) node[below] {\small$\x$};

  % Oznaczenie cięcia
  \node[below, bgreen] at (7,-0.3) {\tiny cięcie $b_1$};
  \node[below, bamber] at (5,-0.5) {\tiny cięcie $b_2$};
\end{tikzpicture}
\caption{Przykład ewolucji wysokości trzech bambusów
($h_1=0{,}7$, $h_2=0{,}2$, $h_3=0{,}1$) według strategii \RM{}.
Czerwona linia przerywana oznacza osiągnięty makespan.}
\label{fig:example}
\end{figure}

% ======================================================================
\section{Znane strategie heurystyczne}
% ======================================================================

\subsection{Strategia Reduce-Max}

\begin{definicja}[\RM]
W każdym dniu ogrodnik przycina bambus o \emph{największej aktualnej wysokości}.
Remisy rozstrzygane są dowolnie.
\end{definicja}

\RM{} jest najbardziej naturalną heurystyką i~zachowuje się bardzo dobrze
w~praktyce. Wyniki eksperymentalne sugerują, że jej makespan może być
ograniczony przez~$2$~\cite{demidio2019}, jednak najlepsze znane ograniczenie
górne w~scenariuszu adwersarialnym wynosi
$1 + \mathcal{H}_{n-1} = \Theta(\log n)$, gdzie $\mathcal{H}_{n-1}$
to $(n-1)$-ta liczba harmoniczna~\cite{bodlaender2012}.

Kluczowym wynikiem artykułu Bilò i~wsp.\ jest:

\begin{twierdzenie}[\cite{bilo2022}, Twierdzenie~1]\label{thm:rm}
\RM{} gwarantuje makespan co najwyżej $9$.
\end{twierdzenie}

Jest to \textbf{pierwsze stałe ograniczenie górne} dla tej strategii
w~scenariuszu nieadwersarialnym, wykazując rozdział między scenariuszem
statycznym a~adwersarialnym.

\subsubsection{Szkic dowodu}

Podzielmy bambusy na \emph{poziomy}: bambus $b_i$ jest na poziomie~$j\geq 1$,
jeśli
\[
  \frac{1}{2^j} \leq h_i < \frac{1}{2^{j-1}}.
\]
Niech $L_j$ oznacza zbiór bambusów poziomu $j$, a~$\sigma(j)$~– maksymalną
wysokość kiedykolwiek osiągniętą przez dowolny bambus poziomu $\geq j$
(z~$\sigma(K+1) = 0$). Kluczowa rekurencja to:

\begin{equation}\label{eq:recurrence}
  \sigma(j) \;\leq\; \max\!\bigl\{3,\,\sigma(j+1)\bigr\}
             \;+\; 3\sum_{k=1}^{j} \frac{|L_k|}{2^j}.
\end{equation}

\textbf{Argument o powtarzanych cięciach.} Rozważmy okno $[d_0+1, d_1]$,
w~którym \RM{} tnie wyłącznie bambusy poziomu $\leq j$ o wysokości $\geq 3$.
Cięcia dzielimy na:
\begin{itemize}[noitemsep]
  \item \emph{pierwsze cięcia} – bambus cięty po raz pierwszy w oknie;
        jest ich co najwyżej $\sum_{k=1}^{j}|L_k|$;
  \item \emph{powtarzane cięcia} – bambus cięty ponownie; każde usuwa
        co najmniej $3$ jednostki wzrostu, więc stanowią co najwyżej
        $\tfrac{1}{3}$ wszystkich dni okna.
\end{itemize}
Stąd $d_1 - d_0 \leq \tfrac{3}{2}\sum_{k=1}^{j}|L_k|$.

Sumując wkłady bambusów i korzystając z $\sum h_i = 1$:
\[
  \sigma(1) \;\leq\; 3 + 3\sum_{j=1}^{K}\sum_{k=1}^{j}\frac{|L_k|}{2^j}
             \;\leq\; 3 + 6\sum_{i=1}^{n} h_i \;=\; 9. \qquad \square
\]

\subsection{Strategia Reduce-Fastest(x)}

\begin{definicja}[\RF$(x)$]
W każdym dniu ogrodnik przycina bambus o \emph{najwyższym współczynniku wzrostu}
spośród tych, które osiągnęły wysokość co najmniej~$x$. Jeśli żaden bambus
nie przekroczył progu~$x$, strategia może nie wykonać cięcia; w~naszej
implementacji stosujemy praktyczny wariant z~awaryjnym wyborem \RM{},
żeby symulacja zawsze wykonywała jedno sensowne działanie dziennie.
\end{definicja}

\begin{twierdzenie}[\cite{bilo2022}, Twierdzenie~2]\label{thm:rf}
Dla stałego $x > 1$, makespan \RF$(x)$ jest ograniczony przez:
\[
  \makespan \;\leq\; \max\!\left\{
    x + \frac{x^2}{4(x-1)},\quad
    \frac{1}{2} + x + \frac{x^2}{4\!\left(x - \tfrac{1}{2}\right)}
  \right\}.
\]
\end{twierdzenie}

\begin{wniosek}\label{cor:golden}
Dla $x = 1 + \tfrac{1}{\sqrt{5}}$ (optymalny wybór) makespan \RF$(x)$ wynosi
co najwyżej $1 + \varphi = \tfrac{3+\sqrt{5}}{2} \approx 2{,}618$,
gdzie $\varphi = \tfrac{1+\sqrt{5}}{2}$ to złota liczba.

Dla $x = 2$ (poprzednio uważany za najlepszy) uzyskujemy poprawione
ograniczenie $\tfrac{19}{6} \approx 3{,}17$, lepsze od wcześniejszego
wyniku~$4$.
\end{wniosek}

\subsubsection{Kluczowe idee dowodu}

Niech $[0, T]$ oznacza interwał, w~którym bambus $b_i$ nie jest cięty
(od momentu gdy osiągnął wysokość $\geq x$). Definiujemy:
\begin{itemize}[noitemsep]
  \item \emph{objętość} $V = T$ – całkowity wzrost ogrodu w~$[0,T-1]$;
  \item $N$ – liczba różnych bambusów ciętych w~$[0,T-1]$;
  \item $V'$ – objętość usunięta przez powtarzane cięcia.
\end{itemize}

Ponieważ każde powtarzane cięcie usuwa co najmniej $x$ jednostek:
$V' \leq T(1-h_i) - N^2 h_i$. Z warunkiem $h_j \geq h_i$ dla każdego
ciętego~$b_j$, wyprowadza się:
\[
  T \;\leq\; \frac{x^2}{4 h_i (h_i + x - 1)}.
\]
Makespan wynosi $M \leq x + h_i + T h_i$, co po optymalizacji po $h_i$
daje wzór z~\cref{thm:rf}.

\subsection{Porównanie heurystyk}

\begin{table}[htbp]
\centering
\caption{Zestawienie znanych granic makespanu dla omawianych strategii.}
\label{tab:comparison}
\small
\begin{tabularx}{\textwidth}{@{}lYYY@{}}
\toprule
\textbf{Strategia} & \textbf{Ograniczenie górne} & \textbf{Ograniczenie dolne}
  & \textbf{Złożoność oracle} \\
\midrule
\RM{} (adwersarialny) & $\Theta(\log n)$, ciasne & $\Theta(\log n)$
  & $O(\log^2 n)$ lub $O(\log n)$ amortyzowane \\
\RM{} (statyczny)     & $\leq 9$ \textbf{[nowy]} & $\geq 1$
  & $O(\log^2 n)$ \\
\RF$(x=2)$            & $\leq 19/6$ \textbf{[poprawa]} & brak znanych
  & $O(\log n)$ \\
\RF$(x=1+1/\sqrt{5})$ & $\leq (3+\sqrt{5})/2 \approx 2{,}62$ \textbf{[nowy]}
  & brak znanych & $O(\log n)$ \\
Optymalny (Pinwheel) & $\leq 2$, asymptotycznie ciasne & $\geq 2-\varepsilon$
  & $O(\log n)$ \\
\bottomrule
\end{tabularx}
\end{table}

% ======================================================================
\section{Powiązanie z problemem Pinwheel Scheduling}\label{sec:pinwheel}
% ======================================================================

\begin{definicja}[Pinwheel Scheduling, PS]
Dane jest $n$~zadań z~całkowitymi częstotliwościami $f_i \geq 1$.
Szukamy wieczystego harmonogramu $S$ takiego, że każde zadanie~$i$
pojawia się co najmniej raz w~każdym spójnym podciągu długości~$f_i$.
\emph{Gęstość} instancji definiujemy jako $\rho = \sum_{i=1}^{n} \tfrac{1}{f_i}$.
\end{definicja}

\begin{twierdzenie}[Warunek konieczny i dostateczny, \cite{holte1989}]
Instancja PS z~$\rho > 1$ nie posiada dopuszczalnego harmonogramu.
Jeśli $\rho \leq 1$ i~wszystkie $f_i$ są potęgami~$2$, dopuszczalny
harmonogram zawsze istnieje.
\end{twierdzenie}

\subsection{Redukcja BGT $\to$ PS dla makespanu 2}

Aby uzyskać makespan $\leq 2$, chcemy ciąć bambus $b_i$ co najwyżej
$\lfloor 2/h_i \rfloor$ dni. Definiujemy:
\[
  f_i \;=\; 2^{\lfloor \log_2(2/h_i) \rfloor}
  \quad\text{(zaokrąglenie $\lfloor 2/h_i \rfloor$ w~dół do potęgi 2)}.
\]
Gęstość wynikowej instancji PS:
\[
  \rho \;=\; \sum_{i=1}^{n} \frac{1}{f_i}
       \;\leq\; \sum_{i=1}^{n} h_i
       \;=\; 1.
\]
Ponieważ $f_i$ są potęgami 2, istnieje dopuszczalny harmonogram PS,
który bezpośrednio daje harmonogram BGT z~makespanem $\leq 2$.

\textbf{Problem praktyczny.} Wszystkie dopuszczalne harmonogramy instancji PS
mogą mieć długość $\Omega\!\bigl(\prod_{i=1}^{n} f_i\bigr)$~\cite{holte1989},
co czyni naiwne podejście wykładniczym. Rozwiązaniem jest \emph{Trimming Oracle}
(patrz \cref{sec:oracles}).

Relacja między BGT i~Pinwheel Scheduling wpisuje się w~szerszy nurt badań
nad harmonogramami okresowymi. Klasy dopuszczalnych instancji Pinwheel były
rozszerzane między innymi w~\cite{chan1993}, natomiast dla dyskretnego BGT
rozwijano także niezależne podejścia aproksymacyjne, np.~\cite{vee2020}.

% ======================================================================
\section{Trimming Oracles – struktury danych}\label{sec:oracles}
% ======================================================================

\begin{definicja}[Trimming Oracle]
Trimming Oracle to struktura danych utrzymująca stan ogrodu i odpowiadająca
na zapytanie \emph{„Który bambus ciąć dzisiaj?"} w~czasie sub-liniowym,
używając przestrzeni liniowej.
\end{definicja}

Poniżej opisujemy trzy wyrocznie (\emph{oracles}) z~artykułu~\cite{bilo2022}.

\subsection{Oracle dla Reduce-Fastest(x)}

\textbf{Idea.} Dla każdego bambusa $b_i$ utrzymujemy dzień $d_i$,
w~którym $b_i$ osiągnie wysokość $\geq x$. Zapytanie zwraca $b_i$
z~minimalnym indeksem $i$ spośród tych, dla których $d_i \leq D$.

\textbf{Struktura danych.} Używamy \emph{drzewa wyszukiwania priorytetowego}
(priority search tree, PST)~\cite{mccreight1985} utrzymującego punkty
$(d_i, i)$ z~operacjami:
\begin{itemize}[noitemsep]
  \item $\textsc{MinYInXRange}(T, x_0)$: minimum $i$ wśród $(d_i, i)$
        z~$d_i \leq x_0$;
  \item $\textsc{GetX}(T, y)$: współrzędna $d_i$ punktu o~$y = i$.
\end{itemize}
Obie operacje działają w~czasie $O(\log n)$.

\textbf{Problem rosnących indeksów.} Dni $d_i$ i~$D$ rosną bezgranicznie.
Rozwiązanie: dzielimy czas na \emph{interwały} $I_j = [nj, n(j+1)-1]$
i~utrzymujemy dwa drzewa $T_1$, $T_2$ dla bieżącego i~następnego
interwału, mierząc dni od początku interwału.

\begin{algorithm}[htbp]
\caption{Trimming Oracle dla \RF$(x)$ – funkcja Query}
\label{alg:rf}
\begin{algorithmic}[1]
\Require dzień $\delta$ w~bieżącym interwale $I_j$
\Ensure indeks bambusa do przycięcia (lub brak akcji)
\State $i \gets \textsc{MinYInXRange}(T_1, \delta)$
\If{$i$ istnieje}
  \State $\textsc{Update}(T_1,\; \delta + \lceil x/h_i \rceil,\; i)$
  \State $\textsc{Update}(T_2,\; \delta + \lceil x/h_i \rceil - n,\; i)$
\EndIf
\State Zaktualizuj $T_2$ dla bambusa $b_{\delta+1}$
\State $\delta \gets (\delta + 1) \bmod n$
\If{$\delta = 0$}
  \State Zamień $T_1$ i~$T_2$
\EndIf
\State \Return „Przytnij $b_i$" (lub „brak akcji")
\end{algorithmic}
\end{algorithm}

\begin{twierdzenie}[\cite{bilo2022}, Twierdzenie~3]\label{thm:oracle-rf}
Istnieje Trimming Oracle implementujący \RF$(x)$ o~przestrzeni $O(n)$,
czasie budowy $O(n \log n)$ i~czasie zapytania $O(\log n)$ w~pesymistycznym
przypadku.
\end{twierdzenie}

\subsection{Oracle dla Reduce-Max}

\textbf{Idea.} Każdy bambus $b_i$ reprezentujemy prostą $\ell_i(d) = h_i d + c_i$
opisującą jego wysokość w~dniu $d$. Zapytanie wymaga znalezienia
$\arg\max_i \ell_i(D)$, co odpowiada obliczeniu \emph{górnej otoczki}
(upper envelope) zbioru prostych.

\textbf{Struktura danych.} Używamy dynamicznej struktury utrzymującej
górną otoczkę~\cite{overmars1981, brodal2002}, obsługując:
\begin{itemize}[noitemsep]
  \item wstawianie i~usuwanie prostej: $O(\log^2 n)$ pesymistycznie
        lub $O(\log n)$ amortyzowanym;
  \item zapytanie $\textsc{Upper}(U, d)$: prosta maksymalna w~punkcie $d$
        w~czasie $O(\log n)$.
\end{itemize}

\begin{twierdzenie}[\cite{bilo2022}, Twierdzenie~4]\label{thm:oracle-rm}
Istnieje Trimming Oracle implementujący \RM{} o~przestrzeni $O(n)$,
czasie budowy $O(n \log n)$, czasie zapytania $O(\log^2 n)$ pesymistycznie
lub $O(\log n)$ amortyzowanym.
\end{twierdzenie}

\subsection{Oracle gwarantujący makespan 2}

\textbf{Idea.} Zaokrąglamy współczynniki $h_i$ do najbliższych mniejszych
potęg dwójki i~konstruujemy wieczyste harmonogramy dla tak zwanych \emph{instancji
regularnych}.

\begin{definicja}[Instancja regularna]
Instancja jest \emph{regularna}, jeśli bambus $b_i$ ma współczynnik wzrostu
$h_i = 2^{-i}$ dla $i = 1, \ldots, n-1$ i~$h_n = h_{n-1} = 2^{-(n-1)}$.
\end{definicja}

Dla instancji regularnych istnieje elegancka reguła: bambus do przycięcia
w~dniu $D$ to $b_i$, gdzie $i$ jest pozycją \emph{najmniej znaczącego bitu
równego 0} w~binarnej reprezentacji $D$. Jeśli wszystkie $n-1$ najmłodszych
bitów wynosi~$1$, tniemy $b_n$.

\begin{twierdzenie}[\cite{bilo2022}, Lemat~1]
Dla instancji regularnych istnieje Trimming Oracle o~przestrzeni $O(n)$,
czasie budowy $O(n)$ i~czasie zapytania $O(1)$ pesymistycznie,
gwarantujący makespan~$1$.
\end{twierdzenie}

Dla ogólnych instancji stosuje się hierarchiczną dekompozycję:
bambusy są iteracyjnie łączone w~\emph{wirtualne bambusy}, tworząc
drzewo $T_O$ o~logarytmicznej głębokości. Każdy węzeł wewnętrzny
zawiera Oracle dla instancji regularnej (skalibrowanej lokalnie).

\begin{twierdzenie}[\cite{bilo2022}, Twierdzenie~5]\label{thm:oracle-opt}
Istnieje Trimming Oracle gwarantujący makespan~$2$, o~przestrzeni $O(n)$,
czasie budowy $O(n \log n)$ i~czasie zapytania $O(\log n)$.
\end{twierdzenie}

% ======================================================================
\section{Implementacja: interaktywna aplikacja webowa}
% ======================================================================

W ramach projektu przygotowaliśmy interaktywną aplikację webową
wizualizującą problem BGT. Aktualna wersja repozytorium jest napisana
w~\textbf{Pythonie} z~użyciem biblioteki \textbf{Streamlit}. Taki wybór
technologii dobrze pasuje do charakteru projektu: najważniejsza jest
czytelna logika symulacji, szybkie eksperymentowanie z~parametrami oraz
proste uruchomienie aplikacji bez osobnego frontendu i~backendu.

Aplikacja nie korzysta z~zewnętrznych bibliotek obliczeniowych. Strategie
\RM{}, \RF$(x)$ oraz harmonogram z~gwarancją makespanu~$2$ są zaimplementowane bezpośrednio
w~kodzie projektu, a~wizualizacje ogrodu i~wykresów generowane są jako
wektorowe rysunki SVG po stronie Pythona.

Implementacja harmonogramu Makespan-2 stosuje redukcję z~\cite{bilo2022}:
normalizuje współczynniki wzrostu, zaokrągla je w~dół do potęg~$1/2$,
dopełnia instancję dyadyczną do sumy~$1$, a~następnie buduje drzewo
wirtualnych bambusów. Każdy węzeł wewnętrzny harmonogramuje swoje dzieci
jako przeskalowaną instancję regularną, zgodnie z~konstrukcją oracle'a
gwarantującego makespan~$2$. Wersja w~aplikacji zachowuje tę konstrukcję
i~gwarancję makespanu, ale dla małych instancji używa prostszego
wyszukiwania operacji merge oraz arytmetycznego wyznaczania kolejnego
dziecka zamiast zoptymalizowanych struktur danych z~dowodu złożoności
$O(n\log n)$ budowy i~$O(\log n)$ zapytania.

\subsection{Architektura}

\begin{figure}[htbp]
\centering
\begin{tikzpicture}[
  box/.style={draw, rounded corners=4pt, minimum width=3cm,
              minimum height=1cm, align=center, font=\small},
  arr/.style={->, >=Stealth, thick}
]
  \node[box, fill=bgreen!15, draw=bgreen] (model) at (0,0)
    {Model danych\\[-2pt]\tiny\textit{Bamboo, growth, height}};
  \node[box, fill=bamber!15, draw=bamber] (logic) at (4,0)
    {Logika symulacji\\[-2pt]\tiny\textit{grow(), cut(), advance\_day()}};
  \node[box, fill=bblue!15, draw=bblue] (svg) at (8,0)
    {Renderowanie SVG\\[-2pt]\tiny\textit{garden + sparkline}};
  \node[box, fill=bred!15, draw=bred] (state) at (2,-2)
    {Stan sesji\\[-2pt]\tiny\textit{st.session\_state}};
  \node[box, fill=gray!15, draw=black!45] (ui) at (6,-2)
    {Interfejs Streamlit\\[-2pt]\tiny\textit{Sandbox / Compare / Game}};

  \draw[arr, bgreen] (model) -- (logic);
  \draw[arr, bamber] (logic) -- (svg);
  \draw[arr] (state) -- (logic);
  \draw[arr] (ui) -- (state);
  \draw[arr] (svg) -- (ui);
\end{tikzpicture}
\caption{Architektura aktualnej aplikacji BGT. Część obliczeniowa jest
wydzielona na czyste funkcje, a~Streamlit odpowiada za stan sesji,
kontrolki i~prezentację.}
\label{fig:arch}
\end{figure}

Najważniejsze elementy kodu:
\begin{itemize}[noitemsep]
  \item \texttt{Bamboo} – niemutowalna struktura opisująca pojedynczy bambus;
  \item \texttt{grow()} i~\texttt{cut()} – podstawowe operacje na ogrodzie;
  \item \texttt{advance\_day()} – wspólna funkcja wykonująca jeden dzień
        symulacji dla dowolnej strategii;
  \item \texttt{oracle\_reduce\_max()}, \texttt{oracle\_reduce\_fastest()}
        oraz \texttt{MakespanTwoOracle} – implementacje reguł i harmonogramu cięć;
  \item \texttt{bamboo\_garden\_svg()} i~\texttt{sparkline\_svg()} –
        generatory wizualizacji w~SVG.
\end{itemize}

\subsection{Tryby działania}

Aplikacja oferuje trzy tryby:

\begin{enumerate}[noitemsep]
  \item \textbf{Sandbox} – użytkownik konfiguruje ogród (liczba bambusów,
        współczynniki wzrostu), wybiera algorytm spośród \RM{}, \RF$(x)$
        (z~regulowanym progiem $x$) i~harmonogramu Makespan-2, obserwuje
        animowaną symulację dzień po dniu.

  \item \textbf{Porównanie} – wszystkie trzy algorytmy działają
        równolegle na tym samym wejściu; na żywo wyświetlane są
        makespany i~historia na wspólnym wykresie.

  \item \textbf{Tryb gry} – użytkownik sam gra rolę ogrodnika,
        klika w~bambusy, by je przyciąć, rywalizując z~botem
        \RM{} o~niższy makespan.
\end{enumerate}

\subsection{Kluczowe decyzje implementacyjne}

\begin{itemize}[noitemsep]
  \item \textbf{Dyskretna symulacja} – każdy krok odpowiada jednemu dniowi:
        bambusy najpierw rosną, następnie oracle wybiera bambus do przycięcia.
  \item \textbf{Niemutowalny model danych} – funkcje \texttt{grow()} i
        \texttt{cut()} zwracają nowe listy bambusów, co upraszcza testowanie
        i~zmniejsza ryzyko przypadkowego nadpisania stanu.
  \item \textbf{Wspólny interfejs strategii} – każda strategia jest funkcją
        przyjmującą listę bambusów i~zwracającą indeks wybranego cięcia.
  \item \textbf{Stan w~\texttt{st.session\_state}} – Streamlit przechowuje
        aktualny ogród, historię makespanu, wybrany tryb i~parametry
        użytkownika.
  \item \textbf{Renderowanie SVG} – grafika jest stylizowana na ilustrację
        z~artykułu: zielone łodygi bambusa, znaczniki wysokości, linie
        pomocnicze i~robot-panda wskazujący ostatnie cięcie.
  \item \textbf{Testowalna warstwa obliczeniowa} – funkcje symulacyjne są
        niezależne od Streamlita i~są pokryte testami jednostkowymi.
\end{itemize}

\subsection{Uruchomienie i testy}

Projekt można uruchomić lokalnie poleceniami:
\begin{verbatim}
python3.12 -m venv .venv
source .venv/bin/activate
pip install -r requirements.txt
streamlit run app.py
\end{verbatim}

Do weryfikacji logiki symulacji przygotowano testy jednostkowe:
\begin{verbatim}
python -m unittest
\end{verbatim}

Testy sprawdzają między innymi niezmienniczość operacji \texttt{grow()}
i~\texttt{cut()}, wybór bambusa przez \RM{} i~\RF$(x)$, normalizację
współczynników wzrostu oraz poprawne raportowanie makespanu przed cięciem.

\subsection{Wkład własny}

Warstwa teoretyczna raportu opiera się na wynikach znanych z~literatury,
przede wszystkim z~pracy Bilò i~wsp.~\cite{bilo2022}. Nasz wkład projektowy
polega na przełożeniu tych wyników na działające narzędzie eksperymentalne:
\begin{itemize}[noitemsep]
  \item przygotowaliśmy implementację symulacji BGT w~Pythonie;
  \item zaimplementowaliśmy trzy reguły wyboru cięcia i~wspólny mechanizm
        kroku symulacji;
  \item zbudowaliśmy trzy tryby interakcji: Sandbox, Compare oraz Game;
  \item dodaliśmy wizualizacje SVG pokazujące aktualne wysokości bambusów,
        ostatnie cięcie i~historię makespanu;
  \item przygotowaliśmy testy jednostkowe sprawdzające kluczowe operacje
        modelu.
\end{itemize}

Dzięki temu raport nie kończy się wyłącznie na omówieniu twierdzeń:
pokazuje również, jak abstrakcyjne strategie zachowują się w~konkretnych,
uruchamialnych instancjach.

% ======================================================================
\section{Analiza złożoności i podsumowanie wyników}
% ======================================================================

\subsection{Zestawienie złożoności Oracle'i}

\begin{table}[htbp]
\centering
\caption{Złożoność czasowa i~przestrzenna Trimming Oracle'i.}
\label{tab:oracles}
\small
\begin{tabularx}{\textwidth}{@{}lYYYY@{}}
\toprule
\textbf{Strategia / Oracle} & \textbf{Makespan} & \textbf{Budowa}
  & \textbf{Zapytanie} & \textbf{Przestrzeń} \\
\midrule
\RM{} naiwny        & $\leq 9$    & $O(n)$        & $O(n)$ & $O(n)$ \\
\RM{} z~otoczką     & $\leq 9$    & $O(n\log n)$  & $O(\log^2 n)$ lub $O(\log n)$ amort. & $O(n)$ \\
\RF$(x)$            & $\leq 2{,}62$ dla opt. $x$ & $O(n\log n)$ & $O(\log n)$ & $O(n)$ \\
Oracle makespan-2   & $\leq 2$    & $O(n\log n)$  & $O(\log n)$ & $O(n)$ \\
\bottomrule
\end{tabularx}
\end{table}

Warto podkreślić różnicę między \emph{strategią} a~\emph{oracle}. Strategia
określa, który bambus powinien zostać przycięty, natomiast oracle odpowiada
za efektywne znalezienie tego bambusa bez przeglądania całego ogrodu.
Nasza aplikacja implementuje reguły wyboru bez zaawansowanych struktur danych,
ponieważ jej celem jest interaktywna demonstracja i~eksperymentowanie na
umiarkowanych instancjach, a~nie obsługa milionów bambusów.

\subsection{Otwarte problemy}

Artykuł pozostawia kilka otwartych pytań o~dużym znaczeniu:

\begin{enumerate}[noitemsep]
  \item \textbf{Hipoteza dla \RM:} Czy makespan \RM{} wynosi co najwyżej~$2$
        dla każdej instancji? Eksperymentalnie odpowiedź jest twierdząca,
        ale dowód algebraiczny pozostaje nieznany.
  \item \textbf{Hipoteza dla \RF$(1)$:} Czy \RF$(1)$ gwarantuje
        makespan~$2$? To odpowiadałoby optymalności przy zerowym narzucie
        na selekcję.
  \item \textbf{Dolne ograniczenie dla \RF$(x)$:} Nie są znane żadne
        dolne ograniczenia inne niż trywialne $\geq 1$.
  \item \textbf{Lepsza praktyczna implementacja oracle'i:} Interesujące
        byłoby porównanie struktur teoretycznych z~prostymi implementacjami
        używającymi kopców, drzew przedziałowych lub struktur bibliotecznych.
\end{enumerate}

% ======================================================================
\section{Wnioski}
% ======================================================================

Problem Bamboo Garden Trimming jest dobrym przykładem sytuacji, w~której
bardzo prosta reguła działania prowadzi do nietrywialnej analizy
kombinatorycznej. Wystarczy jeden robot, jeden ogród i~jedno cięcie dziennie,
żeby pojawiły się pytania o~granice aproksymacji, redukcje do harmonogramowania
wieczystego oraz dynamiczne struktury danych.

Główne wyniki artykułu Bilò i~wsp.~\cite{bilo2022} to:

\begin{itemize}[noitemsep]
  \item pierwsze stałe ograniczenie górne na makespan \RM{}: $\leq 9$,
  \item poprawa ograniczenia dla \RF$(x)$ do $\approx 2{,}62$
        przy optymalnym $x = 1 + 1/\sqrt{5}$,
  \item trzy wydajne Trimming Oracle'i, w~tym Oracle gwarantujący
        makespan~$2$ w~czasie $O(\log n)$ na zapytanie,
  \item rozwiązanie otwartego problemu z~\cite{gasieniec2017}
        dotyczącego implementacji \RM{} w~czasie $o(n)$.
\end{itemize}

Zaimplementowana przez nas aplikacja webowa pozwala na interaktywne
badanie powyższych strategii i~wizualną weryfikację ich zachowania
na wybranych instancjach. Tryb porównania pokazuje, że różnice między
\RM{}, \RF$(x)$ i~harmonogram Makespan-2 nie są jedynie abstrakcyjne: dla tych samych
współczynników wzrostu algorytmy mogą generować inną historię makespanu
i~inne decyzje cięcia. Projekt stanowi więc praktyczne uzupełnienie analizy
teoretycznej: pozwala zobaczyć, jak twierdzenia przekładają się na dynamikę
konkretnego ogrodu.

% ======================================================================
% BIBLIOGRAFIA
% ======================================================================
\newpage
\begin{thebibliography}{99}

\bibitem{bilo2022}
D.~Bilò, L.~Gualà, S.~Leucci, G.~Proietti, G.~Scornavacca,
\textit{Cutting bamboo down to size},
Theoretical Computer Science \textbf{909} (2022), ss.~54--67.
\url{https://doi.org/10.1016/j.tcs.2022.01.027}

\bibitem{gasieniec2017}
L.~G\k{a}sieniec, R.~Klasing, C.~Levcopoulos, A.~Lingas, J.~Min,
T.~Radzik,
\textit{Bamboo garden trimming problem (perpetual maintenance of machines
with different attendance urgency factors)},
w: SOFSEM 2017, Lecture Notes in Computer Science,
Springer 2017, ss.~229--240.

\bibitem{holte1989}
R.~Holte, A.~Mok, L.~Rosier, I.~Tulchinsky, D.~Varvel,
\textit{The pinwheel: a real-time scheduling problem},
Proceedings of the 22nd Annual Hawaii International Conference
on System Sciences, 1989, ss.~693--702.

\bibitem{bodlaender2012}
M.H.L.~Bodlaender, C.A.J.~Hurkens, V.J.J.~Kusters, F.~Staals,
G.J.~Woeginger, H.~Zantema,
\textit{Cinderella versus the wicked stepmother},
w: TCS 2012, Lecture Notes in Computer Science, t.~7604,
Springer 2012, ss.~57--71.

\bibitem{demidio2019}
M.~D'Emidio, G.D.~Stefano, A.~Navarra,
\textit{Bamboo garden trimming problem: priority schedulings},
Algorithms \textbf{12}(4) (2019), art.~74.
\url{https://doi.org/10.3390/a12040074}

\bibitem{alshamrani2015}
S.S.~Alshamrani, D.R.~Kowalski, L.A.~G\k{a}sieniec,
\textit{Efficient discovery of malicious symptoms in clouds via
monitoring virtual machines},
IEEE CIT/IUCC/DASC/PICom 2015, ss.~1703--1710.

\bibitem{bender2020}
M.A.~Bender, R.~Das, M.~Farach-Colton, R.~Johnson, W.~Kuszmaul,
\textit{Flushing without cascades},
Proceedings of SODA 2020, ss.~650--669.

\bibitem{goldwasser2010}
M.H.~Goldwasser,
\textit{A survey of buffer management policies for packet switches},
SIGACT News \textbf{41}(1) (2010), ss.~100--128.

\bibitem{mccreight1985}
E.M.~McCreight,
\textit{Priority search trees},
SIAM Journal on Computing \textbf{14}(2) (1985), ss.~257--276.

\bibitem{overmars1981}
M.H.~Overmars, J.~van Leeuwen,
\textit{Maintenance of configurations in the plane},
Journal of Computer and System Sciences \textbf{23}(2) (1981),
ss.~166--204.

\bibitem{brodal2002}
G.S.~Brodal, R.~Jacob,
\textit{Dynamic planar convex hull},
Proceedings of FOCS 2002, ss.~617--626.

\bibitem{chan1993}
M.Y.~Chan, F.Y.L.~Chin,
\textit{Schedulers for larger classes of pinwheel instances},
Algorithmica \textbf{9}(5) (1993), ss.~425--462.

\bibitem{vee2020}
M.~van Ee,
\textit{A 12/7-approximation algorithm for the discrete bamboo
garden trimming problem},
CoRR, arXiv:2004.11731, 2020.

\end{thebibliography}

\end{document}
